% !TEX program = xelatex
\documentclass[a4paper,12pt]{article}
\usepackage[UTF8]{ctex}
\usepackage{amsmath,amssymb,amsfonts}
\usepackage{graphicx}
\usepackage{float}
\usepackage{geometry}
\usepackage{caption}
\usepackage{subcaption}
\usepackage{cite}
\usepackage{booktabs}
\usepackage{multirow}
\usepackage{array}
\usepackage{algorithm}
\usepackage{algorithmic}
\usepackage{xcolor}
\usepackage{hyperref}
\usepackage{listings}
\usepackage{longtable}

\geometry{left=2.5cm,right=2.5cm,top=2.5cm,bottom=2.5cm}

\newfontfamily\codefont{Consolas}[Scale=0.8]

\lstset{
  basicstyle=\codefont\small,
  breaklines=true,
  frame=single,
  numbers=left,
  numberstyle=\tiny,
  keywordstyle=\color{blue},
  commentstyle=\color{gray},
  stringstyle=\color{red},
}

\renewcommand{\algorithmicrequire}{\textbf{输入:}}
\renewcommand{\algorithmicensure}{\textbf{输出:}}

\title{\textbf{``浙超''联赛分组方案的数学建模与优化}}
\author{}
\date{\today}

\begin{document}

\maketitle

\begin{abstract}
本文针对2025年浙江省市县联赛（``浙超''）的分组问题，建立了系统的数学模型与优化方案。64支参赛队伍（11个市级队和53个县级队）需分为16个小组，每组4队，并满足严格的行政隶属约束。

针对分组方案生成，采用约束规划（CP-SAT）与模拟退火相结合的方法，生成了6个全部可行的分组方案，所有方案均满足硬约束且软约束违反次数为0。针对抽签方案，设计了分层随机抽签、合法性验签与必要重抽流程，并通过10000次蒙特卡洛模拟验证：硬约束可行率为100\%，平均软约束违反次数为$1.53\pm1.40$。针对比赛地点选择，采用候选场馆筛选与贪心分配算法，确定8个比赛城市，平均旅行距离124.5km，变异系数0.636。针对赛制建议，基于Bradley-Terry实力模型和蒙特卡洛模拟，对比了当前赛制、瑞士轮制和近似双败淘汰制，发现近似双败淘汰制的Gini系数最低（0.4167），且实力分数与夺冠概率的Spearman相关最高（0.9807），公平性表现较优。

本文的方法具有良好的可扩展性和可解释性，可为类似体育赛事分组问题提供参考。
\end{abstract}

\newpage

\section{问题重述}

2025年浙江省举办``浙超''市县联赛，参赛队伍包括11个市级队和53个县级队，共64支队伍。赛事要求将64支队伍分为16个小组（每组4队），先进行小组单循环赛，再进行32强淘汰赛决出冠军。

分组需满足以下约束：
\begin{enumerate}
    \item \textbf{硬约束1}：11个市级队必须分配在不同小组（每组最多1个市级队）；
    \item \textbf{硬约束2}：市级队不能与该市代管的县级队同组；
    \item \textbf{软约束}：同一市代管的县级队尽量分配在不同小组。
\end{enumerate}

本文需要完成四个子任务：
\begin{enumerate}
    \item 生成符合约束的分组方案；
    \item 设计公平透明的抽签方案；
    \item 选择8个比赛地点；
    \item 基于定量分析提出赛制改进建议。
\end{enumerate}

\section{模型假设与符号说明}

\subsection{基本假设}
\begin{itemize}
    \item 以GDP作为队伍实力的代理指标（经济强市通常体育资源更丰富）；
    \item 比赛结果采用Bradley-Terry模型描述胜负概率；
    \item 地理距离采用Haversine公式计算球面距离；
    \item 抽签过程在预先公布的允许集合内随机进行，各队在满足约束的可行位置中近似均衡。
\end{itemize}

\subsection{符号说明}
\begin{table}[H]
\centering
\begin{tabular}{ll}
\toprule
符号 & 含义 \\
\midrule
$T$ & 全部64支队伍的集合 \\
$V$ & 11个市级队伍集合 \\
$C$ & 53个县级队伍集合 \\
$G$ & 16个小组，每组容量4 \\
$S_i$ & 第$i$市代管的县级队伍集合 \\
$s_t$ & 队伍$t$的实力分数 \\
$p_{ij}$ & 队伍$i$战胜队伍$j$的概率 \\
$d(i,j)$ & 两地球间的球面距离(km) \\
\bottomrule
\end{tabular}
\end{table}

\section{分组方案生成（Task 1）}

\subsection{问题分析}
分组问题的核心是在满足硬约束的前提下，优化软约束。硬约束包括：每组恰好4队、每组最多1个市级队、市级队不与代管县同组。软约束要求同市县级队尽量不同组。

\subsection{求解算法}

\subsubsection{约束规划（CP-SAT）}
首先使用Google OR-Tools的CP-SAT求解器生成满足硬约束的可行解。决策变量定义为：
\[
x_{ij} = \begin{cases} 1, & \text{队伍}i\text{分到组}j \\ 0, & \text{否则} \end{cases}
\]

约束条件：
\begin{align}
& \sum_{j=1}^{16} x_{ij} = 1, \quad \forall i \in T \quad \text{(每队恰好分到1组)} \\
& \sum_{i=1}^{64} x_{ij} = 4, \quad \forall j \in G \quad \text{(每组恰好4队)} \\
& \sum_{i \in V} x_{ij} \leq 1, \quad \forall j \in G \quad \text{(每组最多1个市级队)} \\
& x_{vj} + x_{cj} \leq 1, \quad \forall v \in V, c \in S_v, j \in G \quad \text{(市-县不同组)}
\end{align}

\subsubsection{模拟退火优化}
在CP-SAT可行解的基础上，使用模拟退火算法优化软约束。优化目标为最小化同市县级队被分到同组的对数。

参数设置：初始温度$T_0=100$，冷却系数$\alpha=0.995$，终止温度$T_f=1$，最大迭代次数根据问题规模调整。邻域操作为随机交换两个组中的县级队（保持硬约束不变）。

\subsection{结果分析}
共生成6个分组方案，其中方案1和方案2由CP-SAT生成，方案3--6由贪心+随机化生成，所有方案均经模拟退火优化。6个方案全部满足硬约束，软约束违反次数均为0。

最优方案（方案1）的分组详情如表\ref{tab:solution1}所示。这里的最优准则为：先最小化软约束违反次数，再最小化组间GDP变异系数。

\begin{table}[H]
\centering
\caption{最优分组方案（方案1，CP-SAT+SA）}
\label{tab:solution1}
\begin{tabular}{ccl}
\toprule
组号 & GDP总和(亿元) & 队伍 \\
\midrule
1 & 2480 & 桐庐, 东阳, 苍南, 象山 \\
2 & 3435 & 衢州, 龙泉, 德清, 兰溪 \\
3 & 23420 & 杭州, 嵊泗, 嘉善, 新昌 \\
4 & 2500 & 松阳, 嵊州, 永康, 临海 \\
5 & 8537 & 台州, 乐清, 武义, 景宁 \\
6 & 17794 & 宁波, 开化, 天台, 青田 \\
7 & 7062 & 金华, 岱山, 云和, 平阳 \\
8 & 3110 & 遂昌, 诸暨, 常山, 瑞安 \\
9 & 8420 & 绍兴, 文成, 三门, 浦江 \\
10 & 5390 & 海宁, 永嘉, 余姚, 义乌 \\
11 & 8562 & 嘉兴, 仙居, 长兴, 江山 \\
12 & 2110 & 宁海, 磐安, 缙云, 玉环 \\
13 & 3542 & 丽水, 泰顺, 安吉, 平湖 \\
14 & 11810 & 温州, 建德, 庆元, 慈溪 \\
15 & 4803 & 舟山, 温岭, 桐乡, 龙港 \\
16 & 5195 & 湖州, 淳安, 海盐, 龙游 \\
\bottomrule
\end{tabular}
\end{table}

该方案GDP均值为7386亿元，变异系数为0.778，表明组间经济实力存在一定差异（主要由杭州市、宁波市GDP较高导致），但在6个可行方案中均衡性较优。

各方案的GDP均衡性和软约束违反情况如图\ref{fig:gdp}和图\ref{fig:soft}所示。

\begin{figure}[H]
\centering
\includegraphics[width=0.8\textwidth]{../output/task1_gdp_balance.png}
\caption{各方案组间GDP均衡性对比}
\label{fig:gdp}
\end{figure}

\begin{figure}[H]
\centering
\includegraphics[width=0.6\textwidth]{../output/task1_soft_violations.png}
\caption{各方案软约束违反次数对比}
\label{fig:soft}
\end{figure}

分组方案的综合评价雷达图如图\ref{fig:radar}所示。

\begin{figure}[H]
\centering
\includegraphics[width=0.7\textwidth]{../output/task1_radar.png}
\caption{分组方案综合评价雷达图}
\label{fig:radar}
\end{figure}

最优方案的分组热力图如图\ref{fig:heatmap}所示。

\begin{figure}[H]
\centering
\includegraphics[width=0.8\textwidth]{../output/task1_best_heatmap.png}
\caption{最优方案分组热力图}
\label{fig:heatmap}
\end{figure}

\section{抽签方案设计（Task 2）}

\subsection{抽签流程}
设计分层随机抽签方案，确保抽签过程公平透明：

\textbf{第一层：市级队抽签}
11个市级队通过随机抽签确定各自所在小组。由于有16个组但只有11个市级队，市级队被分配到编号1--11的组中。

\textbf{第二层：县级队抽签}
每个市的县级队依次进行抽签。每次抽签时，允许的组为：
\[
A_c = \{g\in G:\ |g|<4,\ g\ne g_{\mathrm{city}}(c)\}
\]
其中$A_c$表示县级队$c$的硬允许组，$g_{\mathrm{city}}(c)$表示其所属市级队所在小组。

在硬允许组中，优先选择尚未被同市其他县级队占据的组；若该优先集合为空，则在硬允许组中随机选择。每轮抽签后进行合法性验签；若局部选择导致后续无可行位置，则更换随机种子重抽，并以随机回溯作为兜底，保证最终结果满足全部硬约束。

\subsection{概率分析与公平性验证}

采用蒙特卡洛方法模拟10000次抽签过程，统计以下指标：

\begin{itemize}
    \item \textbf{硬约束可行率}：100\%的抽签结果满足全部硬约束；
    \item \textbf{平均软约束违反}：$1.53 \pm 1.40$次；
    \item \textbf{均匀性}：市级队在11个市级席位中近似均匀，县级队在满足硬约束的可行组集合内近似均衡。
\end{itemize}

违反次数分布：
\begin{itemize}
    \item 0次违反：约25.0\%
    \item 1次违反：约34.0\%
    \item 2次违反：约19.5\%
    \item 3次及以上：约21.5\%
\end{itemize}

抽签公平性热力图如图\ref{fig:fairness}所示。

\begin{figure}[H]
\centering
\includegraphics[width=0.8\textwidth]{../output/task2_fairness.png}
\caption{抽签公平性：各队落入各组的概率分布}
\label{fig:fairness}
\end{figure}

\section{比赛地点选择（Task 3）}

\subsection{选址方法}
采用两阶段方法选择比赛地点：

\textbf{第一阶段：候选场馆筛选}
以11个市级城市作为初始候选场馆，枚举选出8个比赛地点的组合。对每个组合，要求每个地点承办2个小组，并计算其最小化旅行距离后的分配效果。

\textbf{第二阶段：贪心优化分配}
将16个小组分配到8个比赛地点（每个地点恰好承办2个组），目标是最小化总旅行距离。采用贪心策略：每次选择能使总距离增量最小的(组, 地点)对进行分配，并在所有8地点候选组合中选择总距离最小者。

旅行距离采用Haversine公式计算：
\[
d(i,j) = 2R \cdot \arcsin\left(\sqrt{\sin^2\frac{\Delta\varphi}{2} + \cos\varphi_1 \cos\varphi_2 \sin^2\frac{\Delta\lambda}{2}}\right)
\]
其中$R=6371$km为地球半径。

\subsection{选址结果}
最终确定8个比赛城市及其承办的小组如表\ref{tab:venues}所示。

\begin{table}[H]
\centering
\caption{比赛地点及承办小组}
\label{tab:venues}
\begin{tabular}{cll}
\toprule
比赛地 & 承办组 & 各组平均距离(km) \\
\midrule
杭州 & 组16(112km), 组11(140km) & 126.0 \\
嘉兴 & 组3(100km), 组13(193km) & 146.5 \\
金华 & 组8(110km), 组1(138km) & 124.0 \\
丽水 & 组4(94km), 组7(133km) & 113.5 \\
衢州 & 组2(94km), 组6(176km) & 135.0 \\
绍兴 & 组10(103km), 组9(117km) & 110.0 \\
台州 & 组12(92km), 组15(144km) & 118.0 \\
温州 & 组5(92km), 组14(155km) & 123.5 \\
\bottomrule
\end{tabular}
\end{table}

整体统计：平均距离124.5km，标准差79.1km，变异系数0.636。各队到比赛地的距离分布如图\ref{fig:distance}所示。

\begin{figure}[H]
\centering
\includegraphics[width=0.8\textwidth]{../output/task3_distance.png}
\caption{各队到比赛地距离分布}
\label{fig:distance}
\end{figure}

比赛地点的地理分布可通过交互式地图（\texttt{output/task3\_venue\_map.html}）查看，其中包含各队位置、比赛地点和连线。

\section{赛制建议（Task 4）}

\subsection{实力模型}
基于GDP数据构建队伍实力分数。采用对数归一化处理GDP数据，并加入随机噪声以反映不确定性：
\[
s_t = 0.3 + 0.7 \times \frac{0.7 \cdot \text{norm}(\ln(\text{GDP}_t)) + 0.3 \cdot \varepsilon_t - \min}{\max - \min}
\]
其中$\varepsilon_t \sim U(0,1)$为随机噪声。实力分数归一化到$[0.3, 1.0]$区间。

比赛胜负采用Bradley-Terry模型：
\[
P(i\text{胜}j) = \frac{s_i}{s_i + s_j}
\]

\subsection{三种赛制对比}

\subsubsection{当前赛制}
16组单循环小组赛（每组6场比赛），每组前2名晋级32强淘汰赛。共需$16 \times 6 + 31 = 127$场比赛。

\subsubsection{瑞士轮制}
64队进行6轮瑞士轮配对（按积分排序后相邻配对），取前32名进入淘汰赛。共需$6 \times 32 + 31 = 223$场比赛。

\subsubsection{近似双败淘汰制}
先进行3轮瑞士轮预选取前32名，再进行近似双败淘汰赛（输两场才出局）。比赛场次较多，且在模拟中夺冠概率分布更加均衡。

\subsection{蒙特卡洛模拟结果}
通过5000次蒙特卡洛模拟，得到各队在三种赛制下的夺冠概率。前5名队伍的对比如表\ref{tab:probs}所示。

\begin{table}[H]
\centering
\caption{各队夺冠概率对比（前5名）}
\label{tab:probs}
\begin{tabular}{lcccc}
\toprule
队伍 & 实力分数 & 当前赛制 & 瑞士轮 & 近似双败 \\
\midrule
宁波 & 1.0000 & 0.075 & 0.075 & 0.073 \\
杭州 & 0.8917 & 0.062 & 0.051 & 0.051 \\
温州 & 0.8800 & 0.046 & 0.055 & 0.048 \\
台州 & 0.8453 & 0.051 & 0.037 & 0.044 \\
嘉兴 & 0.8268 & 0.045 & 0.044 & 0.043 \\
\bottomrule
\end{tabular}
\end{table}

Gini系数分析：
\begin{itemize}
    \item 当前赛制：0.4990
    \item 瑞士轮制：0.4847
    \item 近似双败淘汰制：0.4167
\end{itemize}

为避免仅用夺冠概率离散程度评价赛制，进一步引入实力--结果一致性指标。Spearman相关系数衡量队伍实力分数与夺冠概率排序的一致程度，数值越高说明赛制越能保留竞技区分度；实力前16队夺冠概率占比则反映强队优势的集中程度。

\begin{table}[H]
\centering
\caption{赛制公平性与竞技区分度指标}
\label{tab:format_metrics}
\begin{tabular}{lccc}
\toprule
赛制 & Gini系数 & Spearman相关 & 实力前16队夺冠概率占比 \\
\midrule
当前赛制 & 0.4990 & 0.9771 & 0.6072 \\
瑞士轮制 & 0.4847 & 0.9706 & 0.5708 \\
近似双败淘汰制 & 0.4167 & 0.9807 & 0.5284 \\
\bottomrule
\end{tabular}
\end{table}

近似双败淘汰制的Gini系数最低，表明夺冠概率分布最为均匀；同时Spearman相关系数最高，说明其并未明显削弱强队优势与竞技排序。因此，该方案在机会均衡性与竞技区分度之间取得了较好的折中。

弱队（实力排名后20\%）的平均出线率为0.2805，说明当前赛制下弱队仍有一定晋级机会，但强弱差异仍然明显。

夺冠概率对比如图\ref{fig:probs}所示。

\begin{figure}[H]
\centering
\includegraphics[width=0.8\textwidth]{../output/task4_probs.png}
\caption{不同赛制下各队夺冠概率对比}
\label{fig:probs}
\end{figure}

\subsection{建议}
综合以上分析，建议采用``小组循环+近似双败淘汰制''的混合赛制。该方案既保留了小组循环赛的观赏性（每队至少打3场比赛），又通过双败淘汰制思想提高了淘汰赛阶段的机会均衡性（Gini系数最低）。

\section{结论与展望}

\subsection{主要结论}
\begin{enumerate}
    \item 采用CP-SAT约束规划与模拟退火相结合的方法，成功生成6个全部可行的分组方案，所有方案均满足硬约束且软约束违反次数为0。
    \item 设计了分层随机抽签方案，经10000次蒙特卡洛验证，硬约束可行率达100\%，平均软约束违反次数为$1.53\pm1.40$。
    \item 通过候选场馆筛选与贪心优化，确定8个比赛城市，平均旅行距离124.5km，变异系数0.636。
    \item 基于Bradley-Terry模型的赛制分析表明，近似双败淘汰制的机会均衡性最优（Gini系数0.4167），且实力--夺冠概率排序一致性最高（Spearman相关0.9807）。
\end{enumerate}

\subsection{改进方向}
\begin{itemize}
    \item 实力模型仅使用GDP作为代理指标，可引入更多因素（如历史战绩、注册球员数等）；
    \item 比赛地点选择可进一步考虑交通基础设施（高铁站点、机场）的便利性；
    \item 抽签方案可引入权重机制，使县级队的分配更加均衡。
\end{itemize}

\begin{thebibliography}{99}
\bibitem{1} Google OR-Tools, ``CP-SAT Solver,'' \url{https://developers.google.com/optimization/cp/cp_solver}.
\bibitem{2} Bradley, R. A., \& Terry, M. E. (1952). Rank analysis of incomplete block designs. \textit{Biometrika}, 39(3/4), 324-345.
\bibitem{3} 浙江省统计局, ``2024年浙江省各市（县）GDP,'' 国家统计局.
\bibitem{4} 浙江省民政厅, ``浙江省行政区划（2025年）.''
\bibitem{5} Kirkpatrick, S., Gelatt, C. D., \& Vecchi, M. P. (1983). Optimization by simulated annealing. \textit{Science}, 220(4598), 671-680.
\end{thebibliography}

\appendix
\section{支撑材料：核心源代码}
\lstinputlisting[language=Python, caption={分组算法核心代码 (src/grouping.py)}]{../src/grouping.py}
\lstinputlisting[language=Python, caption={抽签方案代码 (src/lottery.py)}]{../src/lottery.py}
\lstinputlisting[language=Python, caption={赛制分析代码 (src/tournament.py)}]{../src/tournament.py}

\end{document}
